% ===========================================================================
%  2 — Definicoes
% ===========================================================================
\chapter{Definições}
\label{cap:definicoes}

Este capítulo fixa o vocabulário do documento. A precisão aqui não é
formalidade: vários termos --- \enquote{família}, \enquote{sessão},
\enquote{documentado} --- têm significado técnico específico neste sistema, e
confundi-los altera a interpretação dos resultados apresentados nos
Capítulos~\ref{cap:ml} e~\ref{cap:rag}.

\section{Domínio: manutenção industrial}

\begin{description}
  \item[Manutenção corretiva] Intervenção realizada após a ocorrência da
    falha, para restabelecer a função do item~\cite{abnt5462}.

  \item[Manutenção preventiva] Intervenção realizada em intervalos
    predeterminados ou segundo critérios prescritos, com o objetivo de reduzir a
    probabilidade de falha, independentemente da condição
    observada~\cite{abnt5462}.

  \item[Manutenção preditiva / baseada em condição (CBM)] Intervenção
    determinada pelo acompanhamento de parâmetros da condição real do
    equipamento --- vibração, temperatura, análise de óleo --- com estimativa do
    instante provável da falha~\cite{iso17359, jardine2006review}.

  \item[Manutenção prescritiva] Extensão da preditiva que, além de detectar e
    prognosticar, recomenda a ação corretiva específica, ordenada e
    fundamentada, integrando conhecimento documental e histórico
    operacional~\cite{ansari2019prima}. É o alvo deste projeto.

  \item[Ativo rotativo] Equipamento cujo princípio de funcionamento envolve
    rotação de eixo --- motor elétrico, bomba, ventilador, redutor. Concentra as
    falhas tratadas neste sistema.

  \item[Bancada de ensaio] Conjunto experimental no qual falhas são
    \textbf{induzidas} de forma controlada para gerar dados rotulados. O
    conjunto de dados deste projeto vem de uma bancada, o que tem consequências
    metodológicas discutidas na Seção~\ref{sec:vazamento}.
\end{description}

\section{Domínio: análise de vibração}

\begin{description}
  \item[Velocidade RMS ($v_{\text{RMS}}$)] Valor eficaz da velocidade de
    vibração, em \si{\milli\meter\per\second}. É o indicador normativo de
    severidade global para máquinas rotativas na faixa de
    \SIrange{10}{1000}{\hertz}~\cite{iso10816, iso20816}.

  \item[Aceleração de pico] Maior valor absoluto de aceleração na janela de
    medição, em $g$. Sensível a impactos, típicos de defeitos localizados em
    rolamentos.

  \item[Curtose ($\kappa$)] Momento estatístico de quarta ordem normalizado.
    Um sinal gaussiano tem $\kappa \approx 3$; valores superiores indicam
    presença de impulsos. Em teoria, é um dos melhores indicadores precoces de
    defeito de rolamento~\cite{randall2011tutorial}; neste conjunto de dados,
    não discrimina (Seção~\ref{sec:assinaturas}).

  \item[Fator de crista ($C_f$)] Razão entre aceleração de pico e aceleração
    RMS. Cresce no início de um defeito localizado e volta a cair quando o
    defeito se generaliza.

  \item[Ordem de rotação] Frequência expressa como múltiplo da frequência de
    rotação do eixo: $\text{ordem} = f / f_{\text{rot}}$, com
    $f_{\text{rot}} = \text{RPM}/60$. Os manuais de manutenção descrevem
    assinaturas em ordens (1$\times$, 2$\times$), não em \si{\hertz} absoluto ---
    razão pela qual a conversão para ordem é uma das \textit{features} derivadas
    do modelo (Seção~\ref{sec:features}).

  \item[BPFO, BPFI, BSF] Frequências características de defeito em rolamentos
    --- pista externa, pista interna e elemento rolante, respectivamente.
    São múltiplos não inteiros da rotação, determinados pela geometria do
    rolamento~\cite{randall2011tutorial}. Sua detecção exige o espectro
    completo, ausente deste conjunto de dados.

  \item[Zonas ISO A/B/C/D] Faixas de severidade da velocidade RMS definidas
    pela norma~\cite{iso10816}: A (máquina recém-comissionada), B (operação
    contínua aceitável), C (insatisfatória --- planejar intervenção) e D
    (inaceitável --- risco de dano). Os limites dependem da classe da máquina
    (I a IV), que é função da potência e da rigidez da fundação.

  \item[Severidade relativa] Razão entre a velocidade RMS medida e a linha de
    base da própria máquina no mesmo regime de rotação. É a leitura mais
    defensável em uma bancada com falhas induzidas, onde a zona absoluta perde
    significado (Seção~\ref{sec:severidade}).
\end{description}

\section{Domínio: dados deste projeto}

\begin{description}
  \item[Evento] Registro único de telemetria: 18 grandezas de sensor,
    identificador, instante de criação e, opcionalmente, a anotação de falha do
    operador. É a unidade atômica de dado do sistema, definida uma única vez em
    \cd{SensorEvent} (\arq{src/core/schemas.py}).

  \item[Rótulo bruto (\cd{raw\_fault})] Texto da coluna \cd{fault} do conjunto
    de dados original, tal como digitado. Contém \num{151} valores distintos,
    incluindo erros de digitação reais (\cd{ddesbalanceado},
    \cd{cockecocked\_adxl\_0}, \cd{mortor\_desligado}) e sufixos de variação
    experimental (\cd{\_2}, \cd{\_carga}, \cd{\_pos\_2}, \cd{new\_}).

  \item[Família canônica] Classe de falha ou estado operacional obtida pela
    canonicalização auditável dos rótulos brutos. São \textbf{17 famílias}: 12
    de falha e 5 de estado. É a única taxonomia usada pelo modelo, pela API e
    pela base documental (Apêndice~\ref{ap:taxonomia}).

  \item[Sessão / campanha de coleta] Conjunto de leituras registradas sob a
    mesma configuração experimental. Neste conjunto de dados, o \emph{rótulo
    bruto identifica a sessão}: \cd{rolamento\_inner},
    \cd{rolamento\_inner\_2} e \cd{new\_rolamento\_inner\_0} são a mesma família
    coletada em três campanhas distintas. Essa coincidência entre sessão e
    rótulo é a causa-raiz do vazamento analisado na
    Seção~\ref{sec:vazamento}.

  \item[Estado operacional] Família cujo \cd{is\_fault} é falso ---
    \cd{normal}, \cd{baseline}, \cd{teste}, \cd{acelerando},
    \cd{motor\_desligado}. Um diagnóstico que aponta estado operacional não
    aciona prescrição alguma.

  \item[Linha de base] Distribuição da velocidade RMS no estado \cd{normal}
    para um regime de rotação específico. Serve de referência para a severidade
    relativa.
\end{description}

\section{Domínio: aprendizado de máquina}

\begin{description}
  \item[\textit{Feature}] Variável de entrada do modelo. O sistema usa 23:
    18 grandezas brutas de sensor e 5 derivadas fisicamente motivadas
    (Seção~\ref{sec:features}).

  \item[Padronização (\textit{StandardScaler})] Transformação
    $z = (x - \mu)/\sigma$ aplicada a cada \textit{feature}. É obrigatória para
    o KNN, cuja distância euclidiana seria dominada pelas grandezas de maior
    magnitude --- neste caso, a rotação, na ordem de \num{1000}.

  \item[KNN] Algoritmo que responde a uma consulta devolvendo os $k$ registros
    mais próximos no espaço de \textit{features}~\cite{cover1967nearest}. Neste
    sistema, $k = 25$ e o índice cobre o histórico completo. É o componente que
    atende literalmente o requisito de \enquote{ocorrências históricas
    semelhantes}.

  \item[LightGBM] Implementação de \textit{gradient boosting} sobre árvores de
    decisão, com amostragem por gradiente unilateral e agrupamento de
    \textit{features} exclusivas~\cite{ke2017lightgbm}. Fornece a classe prevista
    e sua probabilidade.

  \item[Concordância dos vizinhos] Fração dos $k$ vizinhos recuperados que
    pertencem à família prevista pelo classificador. É uma segunda evidência,
    independente da probabilidade, sobre a confiabilidade do diagnóstico.

  \item[\textit{Holdout} por sessão] Protocolo de validação em que sessões
    inteiras de coleta são separadas para teste, nunca leituras individuais.
    Mede a pergunta de produção --- \enquote{o modelo reconhece a falha em uma
    campanha nova?} --- e é a prática recomendada para dados com estrutura de
    agrupamento~\cite{roberts2017crossvalidation}.

  \item[Vazamento de dados (\textit{data leakage})] Situação em que informação
    indisponível no momento da predição real influencia o treino, inflando
    artificialmente a métrica~\cite{kaufman2012leakage, kapoor2023leakage}.
    Neste projeto, o vazamento é de \emph{grupo}: a identidade da campanha de
    coleta está impressa em quase todo o vetor de \textit{features}.

  \item[F1 macro] Média não ponderada do F1 de cada classe. Trata todas as
    classes como igualmente importantes, o que é adequado em conjuntos
    desbalanceados~\cite{sokolova2009systematic}.
\end{description}

\section{Domínio: recuperação e geração}

\begin{description}
  \item[\textit{Embedding}] Representação vetorial densa de um texto, na qual
    proximidade geométrica aproxima similaridade semântica. Este sistema usa
    \cd{gemini-embedding-2} com dimensionalidade 768.

  \item[\textit{Chunk}] Fragmento de documento indexado individualmente
    (aproximadamente \num{1200} caracteres, com sobreposição de 150). Carrega os
    metadados de rastreio --- documento, seção, página, família --- que se
    tornam as citações exibidas ao usuário.

  \item[Base vetorial] Banco especializado em busca por similaridade. Aqui,
    ChromaDB em modo persistente, com índice HNSW~\cite{malkov2020hnsw} e
    métrica de cosseno.

  \item[RAG (\textit{Retrieval-Augmented Generation})] Padrão em que o modelo
    de linguagem recebe, no \textit{prompt}, trechos recuperados de uma base
    externa e responde restrito a eles~\cite{lewis2020rag}. Substitui
    conhecimento paramétrico por conhecimento verificável.

  \item[Contexto fechado] Regime de \textit{prompting} no qual o modelo é
    instruído a responder \textbf{exclusivamente} com base nos trechos
    fornecidos e a declarar explicitamente quando eles não cobrem a pergunta.

  \item[\textit{Guardrail} de cobertura] Verificação determinística, anterior a
    qualquer chamada ao modelo, de que a família de falha possui \textit{chunks}
    indexados. Se não possui, o modelo não é invocado: o sistema devolve o aviso
    de \enquote{não documentado} e a sugestão de registro. É a implementação do
    requisito \req{RF-04}.

  \item[Citação determinística] Referência a documento, seção e página obtida
    dos \emph{metadados} do trecho recuperado, e não do texto gerado pelo
    modelo. Torna a citação alucinada estruturalmente impossível.

  \item[Alucinação] Geração de conteúdo fluente e plausível, mas não sustentado
    pela fonte~\cite{ji2023hallucination}. O Capítulo~\ref{cap:rag} descreve as
    cinco barreiras adotadas contra ela.

  \item[Agente] Modelo de linguagem com autonomia para escolher e invocar
    ferramentas em vários passos até responder. Neste sistema, restringe-se ao
    \textit{endpoint} de conversa; o caminho crítico de diagnóstico permanece
    determinístico (\adrr{ADR-11}).

  \item[Ferramenta (\textit{tool})] Função Python exposta ao agente, com
    assinatura e \textit{docstring} que servem de contrato. São quatro:
    consulta ao histórico, busca documental, diagnóstico e cobertura documental.
\end{description}

\section{Domínio: arquitetura e infraestrutura}

\begin{description}
  \item[Modelo C4] Notação de visualização arquitetural em quatro níveis de
    zoom --- Contexto, Contêiner, Componente e Código ---
    proposta por Simon Brown~\cite{brown2026c4}.

  \item[ADR (\textit{Architecture Decision Record})] Registro curto e imutável
    de uma decisão arquitetural, contendo contexto, decisão, consequências e
    alternativas rejeitadas~\cite{nygard2011adr}.

  \item[Contêiner (C4)] Unidade executável ou de armazenamento que hospeda
    código --- neste sistema: API, interface, \textit{worker}, \textit{broker},
    banco relacional e base vetorial. Não confundir com contêiner Docker, ainda
    que aqui coincidam.

  \item[Infraestrutura como código (IaC)] Declaração da topologia de nuvem em
    arquivo versionado, aplicada por ferramenta. Aqui,
    \arq{.railway/railway.ts}, que é a fonte de verdade da infraestrutura.

  \item[DLQ (\textit{Dead-Letter Queue})] Destino das mensagens que não podem
    ser processadas, acompanhadas do motivo da rejeição. Garante que nenhum
    dado inválido seja descartado em silêncio.

  \item[QoS 1 (MQTT)] Nível de garantia \textit{at-least-once}: a mensagem é
    entregue ao menos uma vez, podendo ser duplicada~\cite{oasis2019mqtt}. Exige
    idempotência no consumidor --- implementada por \cd{session.merge} sobre a
    chave primária do evento.

  \item[Volume persistente] Armazenamento cujo conteúdo sobrevive à
    substituição do contêiner. Necessário para o índice vetorial, que de outro
    modo seria perdido a cada implantação.
\end{description}
